\documentclass[output=paper]{langscibook}
\ChapterDOI{10.5281/zenodo.22880244}
\title{\textit{O poosusu} – The sea poison tree}
\author{Ruth Sima Rigamu and Joyce Maion}
\abstract{}
\begin{document}
\maketitle\label{chap:9}
\begin{figure}[H]
    \includegraphics[width=.65\linewidth]{figures/9_1.jpg}
    \caption{\textit{O poosusu} -- The sea poison tree}
    \label{fig:9_1}
\end{figure}
\begin{xltabular}{\textwidth}{lXX}
    (\ref{ex:text9-1})&\textit{Bea abana huha bona kave to dao raara bona vunuu hata, eori repaa paku bono ogaoga bona vua poosusu.}&When people tie the net that we call \textit{vunuu hata} `bad hair', they make the floats out of sea poison tree fruit.\\
    (\ref{ex:text9-2})&\textit{Momohu na osi maari bona vua poosusu.}&First, they pick sea poison tree fruit.\footnote{\textit{Baringtonia asiatica}}\\
    (\ref{ex:text9-3})&\textit{Eori repaa mosi kopaa bari.}&Then they cut them lengthwise into halves.\\
    (\ref{ex:text9-4})&\textit{Murinae eori repaa toto bono vua komanae.}&After that they remove the seeds from their inside.\\
    (\ref{ex:text9-5})&\textit{Eori repaa nibunibusu pete bona pesun.}&Then they make holes in the fruit husks.\\
    (\ref{ex:text9-6})&\textit{A pesun to nibunibusu riori, eve to vaahana petepete bata riori tea roveana,}&The husks into which they make holes, they hang them one after the other on the seam\\
    (\ref{ex:text9-7})&\textit{ore suguna komana tea vavaarihi na kave.}&from the beginning to the end of net.\\
    (\ref{ex:text9-8})&\textit{A pesun poosusu vai to hana nana tea roveana to pasi agaaga bata nom.}&The sea poison tree husks that are hanging on the seam will be floating.\\
    (\ref{ex:text9-9})&\textit{A vua poosusu to paku riori bona ogaoga.}&The sea poison tree fruit are used to make floats.\\
\end{xltabular}
\begin{figure}[H]
    \captionsetup[subfigure]{oneside,margin={0cm,0cm}}
    \captionsetup{justification=centering}
    \subfloat[]{\includegraphics[trim=0 0 0 50,clip,width=.38\linewidth]{figures/9_2.jpg}\label{fig:9_2}}
    \subfloat[]{\includegraphics[trim=0 20 0 20,clip,width=.24\linewidth]{figures/9_3.jpg}\label{fig:9_3}}
    \subfloat[]{\includegraphics[trim=0 0 0 0,clip,width=.38\linewidth]{figures/9_4.jpg}\label{fig:9_4}}
    \caption{(a) \textit{O purupuru} -- The flower, (b) \textit{O vua} -- The fruit,\\ (c) \textit{O vua kopaa} -- Fruit cut into halves}
\end{figure}
\end{document}